\documentstyle[%


righttag%


,11pt%


,amssymb%


,russian%               remove in Eng.


,izvru%                 Set izven% in Eng.


]{amsart}





%





\theoremstyle{definition}


\theorembodyfont{\rm}


\newtheorem{defnnonum}{\hskip\parindent Определение}


\renewcommand{\thedefnnonum}{}








\newcommand{\rank}{\operatorname{rank}}


\newcommand{\const}{\operatorname{const}}


\newcommand{\up}[2]{\overset{#1}{#2}{}}


\newcommand{\tts}[1]{}





%\hoffset=-15mm%                  For Eng.


\setcounter{page}{23}


\god{2004}


\nomer{11~(510)} %                        Remove in Eng.


%\nomer{Vol.\,48, No.\,11}%               Set in Eng.


%\str{\pageref{begin}--\pageref{end}}%  Set in Eng.


\udk{514.764}





\author{\it М.Л.\,Гаврильченко, В.А.\,Киосак, Й.\,Микеш}


% NB:  ^^ remove in Eng.


\title{Геодезические деформации гиперповерхностей римановых пространств}


\thanks{Поддержано грантом \No\,201/02/0616 грантового агентства 


Республики Чехия.}





\date{}


%\pagestyle{myheadings}%         Set in Eng.


\pagestyle{plain} %              Remove in Eng.


\begin{document}


%\thispagestyle{plain}%          Set in Eng.


\maketitle


\label{begin}








{\bf 1. Введение.}


В теории геодезических отображений римановых пространств


предполагается такое соответствие между их точками, при котором 


{\it каждая} геодезическая линия одного пространства переходит


точно в геодезическую другого


\cite{le}--\cite{ramiga}.





В данной работе рассматриваются такие бесконечно малые


деформации римановых пространств, при которых каждая


геодезическая переходит в геодезическую деформированного


пространства с некоторой контролируемой точностью \cite{ramiga}--\cite{ga89}.


Такой подход, если иметь в виду


прикладные вопросы моделирования, может оказаться даже более


соответствующим реальной физической истории, проистекающей в


гравитационном или электромагнитном поле, при реальном движении


некоторой механической системы и т.\,п.





Отметим, что на сигнатуру метрик римановых пространств $V_n$ не


налагаем ограничения, как принято, например, в 


\cite{pe}--\cite{ramiga}. 


Исследования ведутся локально в классе достаточно гладких


функций. 


\medskip





{\bf 2. Бесконечно малые геодезические деформации римановых пространств.}


Пусть $V_m$ --- риманово пространство, отнесенное к локальным


координатам $(y^1,y^2,\dotsc ,y^m)$, и пусть формулами


\begin{equation}


y^\alpha =y^\alpha (x^1,x^2,\dotsc ,x^n),\quad \rank\bigg\|


\frac{\partial y^\alpha }{\partial x^i}\bigg\| =n<m,


\end{equation}


определяется некоторое изометрически погруженное риманово


подпространство $V_n\subset V_m$. 


Здесь и далее греческие индексы 


$\alpha,\beta,\dotsc$ принимают значения из множества


$\{1,2,\dotsc ,m\}$, латинские 


$i,j,\dotsc$ --- из множества $\{1,2,\dotsc,n\}$.





Если $a_{\alpha \beta }$ и $g_{ij}$ --- метрические тензоры $V_m$


и $V_n$, то \cite{ei}


\begin{equation}


g_{ij}=a_{\alpha \beta }\frac{\partial y^\alpha }{\partial x^i}


\frac{\partial y^\beta }{\partial x^j}.


\end{equation}





Пусть, далее, $\xi^\alpha (x^1,x^2,\dotsc ,x^n)$ --- некоторое


контравариантное векторное поле $V_m$, заданное в точках $V_n$.


Формулы


\begin{equation}


\widetilde{y}^\alpha =y^\alpha (x)+\varepsilon \xi^\alpha (x),


\end{equation}


где $\varepsilon $ --- малый числовой параметр, определяют некоторое 


$\widetilde{V}_n$, которое будем называть 


{\it бесконечно малой деформацией} $V_n$.


Поле $\xi^\alpha (x)$ естественно назвать полем 


{\it смещений} или вектором {\it смещений} \cite{ga89}.





Изучая $\widetilde{V}_n$, будем считать, что величинами порядка 


$\varepsilon^2$ и выше в силу малости $\varepsilon $ или в силу


достаточной точности можно пренебречь.





Таким образом, рассматривается бесконечно малая деформация 1-го


порядка, хотя два последних слова в названии теории будут для


краткости опускаться, как это принято обычно.





В силу сказанного для метрического тензора $\widetilde{V}_n$


\begin{equation}


\widetilde{g}_{ij}=g_{ij}+\varepsilon \delta g_{ij}.


\end{equation}


Найдем тензор $\delta g_{ij}$. Для этого заметим, что формула (3)


дает $\delta y^\alpha =\xi^\alpha (x)$, а 


$$


a_{\alpha \beta }(\widetilde{y})=a_{\alpha \beta }(y+\varepsilon 


\xi)=a_{\alpha \beta }(y)+


\frac{\partial a_{\alpha \beta }}{\partial y^\gamma }


\xi^\gamma \cdot \varepsilon +\dotsb,


$$


так что $\delta a_{\alpha \beta }=


\frac{\partial a_{\alpha \beta }}{\partial y^\gamma }\xi^\gamma$. 





По отношению к преобразованию координат в пространстве $V_n$ все


$y^\alpha$ и $\xi^\alpha $ инвариантны, поэтому их первые


частные производные совпадают с~ковариантными. Заметив все это,


в силу соотношения (2) имеем


\begin{equation}


\delta a_{ij}=


\delta (a_{\alpha \beta }y_{,i}^\alpha y_{,j}^\beta )=


\delta a_{\alpha \beta }y_{,i}^\alpha y_{,j}^\beta +


a_{\alpha \beta }\delta y_{,i}^\alpha y_{,j}^\beta +


a_{\alpha \beta }y_{,i}^\alpha \delta y_{,j}^\beta =


\frac{\partial a_{\alpha \beta }}{\partial y^\gamma }\xi^\gamma


y_{,i}^\alpha y_{,j}^\beta +a_{\alpha \beta }(\xi_{,i}^\alpha y_{,j}^\beta


+y_{,i}^\alpha \xi_{,j}^\beta ).


\end{equation}


Формула (5) справедлива для любой бесконечно малой деформации и


в любой системе координат.





\begin{defnnonum}


Бесконечно малые деформации $V_n$ будем называть 


{\it геодезическими}, если при этих деформациях сохраняются


геодезические линии $V_n$.


\end{defnnonum}





Другими словами, $V_n$ и $\widetilde{V}_n$, отнесенные к общим


координатам $\{x^i\}$, должны допускать геодезическое отображение


друг на друга. Но в таком случае для $V_n$ и $\widetilde{V}_n$


должны выполняться уравнения Леви-Чивита \cite{le}--\cite{ramiga}


\begin{equation}


\widetilde{g}_{ij,k}=2\psi_k\widetilde{g}_{ij}+\psi_i\widetilde{g}


_{jk}+\psi_j\widetilde{g}_{ik},


\end{equation}


причем $\psi_i$ --- градиентный вектор, который определяется инвариантом


$$


\psi=\frac 1{2(n+1)}\ln \bigg| \frac{\widetilde{g}}g\bigg|,


\quad\text{где}\quad


g\overset{\mathrm{def}}{=}{}\det\|g_{ij} \|


\quad\text{и}\quad


\widetilde g\overset{\mathrm{def}}{=}{}\det\|\widetilde g_{ij} \|.


$$


В нашем случае 


$$


2(n+1)\psi =\ln \bigg| \frac{\widetilde{g}}g\bigg| =\ln \bigg| 1+


\frac{\delta g}g\varepsilon \bigg| =\ln \bigg(


1+\frac{\delta g}g\varepsilon \bigg)


=\varepsilon \frac{\delta g}g+\dotsb, 


$$


так что $\psi_i$ в (6) надо заменить на $\varepsilon \psi_i$.





С учетом (4) уравнения Леви-Чивита (6) дают 


\begin{equation}


\delta g_{ij,k}=2\psi_kg_{ij}+\psi_ig_{jk}+\psi_jg_{ik}.


\end{equation}


Наоборот, при выполнении (7) для тензора 


$\widetilde{g}_{ij}=g_{ij}+\varepsilon \delta g_{ij}$ будут выполнены и (6).


Невырожденность $\widetilde{g}_{ij}$ обеспечивается нужным


выбором параметра $\varepsilon $.


Таким образом, условия (7) необходимы и достаточны для того, чтобы


деформация $V_n$ была геодезической.


Симметрический тензор $\delta g_{ij}$


для удобства в дальнейшем будем обозначать через $h_{ij}$.





Таким образом получена 





\begin{thm}[{\series{m}\shape{n}\selectfont \cite{ga72}}]


Для того чтобы риманово пространство $V_n$ допускало бесконечно


малые геодезические деформации, необходимо и достаточно, чтобы в


нем существовал симметрический тензор $h_{ij}$, удовлетворяющий уравнениям


\begin{equation}


h_{ij,k}=2\psi_kg_{ij}+\psi_ig_{jk}+\psi_jg_{ik}


\end{equation}


при некотором градиентном векторе $\psi_i$.


\end{thm}








{\bf 3. Геодезические деформации и геодезические отображения.}


Если в (8) \ $\psi_i=0$, то $h_{ij,k}=0$, что возможно


в двух случаях:


\begin{itemize}


\item[1.] $h_{ij}=Cg_{ij}$, $C=\const$, но тогда (см. (4))


$\widetilde{g}_{ij}=g_{ij}(1+\varepsilon C)$, т.\,е. деформация


сводится к бесконечно малой гомотетии.


\item[2.] $h_{ij}\neq Cg_{ij}$. В этом случае в $V_n$ существует


симметрический ковариантно постоянный тензор.


\end{itemize}


Из (6) следует, что в обоих случаях 


$\widetilde{g}=Cg$, т.\,е. деформация является гомотетически


ареальной. Будем считать эти случаи {\it тривиальными}, нетривиальный случай


характеризуется условием $\psi_i\neq 0$. Но тогда (8) можно


переписать в виде


$(h_{ij}-2\psi g_{ij})_{,k}=\psi_ig_{jk}+\psi_jg_{ik}$, 


т.\,е. в $V_n$ тензор $a_{ij}=h_{ij}-2\psi g_{ij}$ при 


$\lambda_i=\psi_i$ удовлетворяет системе основных уравнений


теории геодезических отображений (см. \cite{si}, с.\,121, формула (8)): 


\begin{equation}


a_{ij,k}=\lambda_ig_{jk}+\lambda_jg_{ik}.


\end{equation}





С другой стороны, для тензора $h_{ij}=a_{ij}+2\lambda g_{ij}$, где


$a_{ij}$ --- решение (9), выполняются уравнения (8),


следовательно метрический тензор $\widetilde{V}_n$ 


имеет представление


$\widetilde{g}_{ij}=(1+2\lambda \varepsilon )g_{ij}+\varepsilon


a_{ij}$. 





Сопоставляя это с теоремой 1, приходим к выводу.





\begin{thm}


Риманово пространство $V_n$ допускает нетривиальные


бесконечно малые геодезические деформации тогда и только тогда,


когда оно допускает и


нетривиальные геодезические отображения.


\end{thm}





Эта теорема впервые была доказана в \cite{siga}, но только для римановых


пространств первого класса (1971~г.).





Из теоремы 2 с учетом теорем Н.С.\,Синюкова и Й.\,Микеша \cite{si},


\cite{mige} вытекает





\begin{thm}


Симметрические, рекуррентные, дважды симметрические, дважды


рекуррентные, $m$-рекуррентные и полусимметрические 


пространства $D^m_n$ непостоянной кривизны не допускают нетривиальных


геодезических деформаций.


\end{thm}





Заметим, что пространства $D^m_n$ были введены в \cite{ve}.





Напротив, для пространств постоянной кривизны, для эквидистантных


пространств такие деформации существуют.


В $E_3$ геодезические деформации допускают поверхности Лиувилля и


только они.


\medskip





{\bf 4. Геодезические деформации подпространств риманова


пространства.} 


Основной задачей теории геодезических деформаций является


всестороннее изучение взаимосвязей и взаимозависимостей риманова


пространства $V_n$, его объемлющего пространства $V_m$ и


изгибающего поля $\xi^\alpha (x)$. Идеальной ситуацией,


очевидно, можно считать случай, когда удается найти поле 


$\xi^\alpha (x)$. Легко получить уравнения для него.





Действительно, подставляя (5) в (7), получим


\begin{multline}


\frac{\partial^2a_{\alpha \beta }}{\partial y^\gamma \partial y^\nu }


y_{,i}^\alpha y_{,j}^\beta \xi^\gamma y_{,k}^\nu +\frac{\partial a_{\alpha


\beta }}{\partial y^\gamma }\big( y_{,i}^\alpha y_{,j}^\beta \xi


_{,k}^\gamma +y_{,i}^\alpha y_{,jk}^\beta \xi^\gamma +y_{,ik}^\alpha


y_{,j}^\beta \xi^\gamma \big) + 


\frac{\partial a_{\alpha \beta }}{\partial y^\gamma }\big(\xi_{,i}^\alpha


y_{,j}^\beta +y_{,i}^\alpha \xi_{,j}^\beta \big) y_{,k}^\gamma+\\


%


+a_{\alpha \beta }\big( \xi_{,ik}^\alpha y_{,j}^\beta +\xi_{,i}^\alpha


y_{,jk}^\beta +y_{,i}^\alpha \xi_{,jk}^\beta +y_{,ik}^\alpha \xi


_{,j}^\beta \big) =


2\psi_kg_{ij}+\psi_ig_{jk}+\psi_jg_{ik}.


\end{multline}


При этом использовано то, что $a_{\alpha \beta }$ инвариантны в


$V_n$. 





Уравнения (10) усложнены тем обстоятельством, что старшая


производная поля смещений $\xi_{,ik}^\alpha $ входит дважды. От


этого можно избавиться, выполнив операцию с такой подстановкой


индексов $i,j,k$: $(i,j,k)+(i,k,j)-(j,k,i)$. Получим


\begin{multline}


a_{\alpha \beta }y_{,i}^\alpha \xi_{,jk}^\beta +a_{\alpha \beta }\xi


_{,i}^\alpha y_{,jk}^\beta +\Gamma_{\gamma \beta \alpha }(y_{,i}^\alpha


y_{,j}^\beta \xi_{,k}^\gamma +y_{,i}^\alpha \xi_{,j}^\beta y_{,k}^\gamma


+\xi_{,i}^\alpha y_{,j}^\beta y_{,k}^\gamma )+  \\


+\xi^\gamma \left( \frac{\partial a_{\alpha \beta }}{\partial y^\gamma }


y_{,i}^\alpha y_{,jk}^\beta +\frac{\partial \Gamma_{\nu \beta \alpha }}{


\partial y^\gamma }y_{,i}^\alpha y_{,j}^\beta y_{,k}^\nu \right) =\psi


_kg_{ij}+\psi_jg_{ik},


\end{multline}


где $\Gamma_{\alpha \beta \gamma }$ -- символы Кристоффеля


пространства $V_m$, составленные для $a_{\alpha \beta }$.


Уравнения (10) и (11) эквивалентны. 


Они представляют необходимые и достаточные условия бесконечно


малых геодезических деформаций $V_n\subset V_m$ с полем смещений 


$\xi^\alpha (x)$.


\medskip





{\bf 5. Основные уравнения геодезических деформаций гиперповерхностей.}


Пусть $m=n+1$, т.\,е. $V_n$ --- гиперповерхность $V_m$. Условие


(1) позволяет векторы $y_{,i}^\alpha $ выбрать в качестве


базиса $V_n$. Будем считать, что 


$\det \| g_{ij}\| \neq 0$, следовательно, нормаль к


гиперповерхности $\eta^\alpha (x)$ неизотропна \cite{ei}, т.\,е.


$$


a_{\alpha \beta }y_{,i}^\alpha \eta^\beta =0,\quad a_{\alpha \beta }\eta


^\alpha \eta^\beta =e,\quad e=\pm 1.


$$


Система векторов $\{y_{,i}^\alpha ,\eta^\alpha \}$ образует


базис $V_m$. Поэтому вектор смещения $\xi^\alpha$ можно


однозначно выразить в следующей форме:


\begin{equation}


\xi^\alpha (x)=\lambda^iy_{,i}^\alpha +\mu \eta^\alpha ,


\end{equation}


где $\lambda^i(x)$ и $\mu(x)$ --- некоторые вектор и инвариант $V_n$. 





Так как $V_n$ --- гиперповерхность $V_{n+1}$, то справедливы


деривационные формулы \cite{ei}


{\it Гаусса:}


$y_{,ij}^{\alpha}=-\Gamma_{\mu \nu }^\alpha y_{,i}^\mu y_{,j}^\nu +e\Omega


_{ij}\eta^\alpha$ 


и {\it Вейнгартена}:


$\eta_{,i}^\alpha=\Gamma_{\mu \nu }^\alpha y_{,i}^\mu \eta^\nu -\Omega


_i^ky_{,k}^\alpha$,


где $\Omega_{ij}$ --- второй основной тензор $V_n$,


$\Omega^k_i\overset{\mathrm{def}}{=}{}\Omega_{ij}g^{jk}$,


$\Gamma_{\mu \nu }^\alpha$ --- символы Кристоффеля пространства $V_m$.





Эти формулы позволяют, используя (12), выразить 


$\xi_{,i}^\alpha $ и $\xi_{,ij}^\alpha$ в том же базисе 


$\{y_{,i}^\alpha ,\eta^\alpha \}$. 


Так, например,


$$


\xi_{,i}^\alpha =(\lambda_{,i}^s-\mu \Omega_i^s)


y_{,s}^\alpha -\lambda^s\Gamma_{\mu \nu }^ay_{,s}^\mu y_{,i}^\nu +(


e\lambda^s\Omega_{si}+\mu_{,i}) \eta^\alpha -\mu \Gamma_{\mu \nu


}^\alpha y_{,i}^\mu \eta^\nu . 


$$


Если все эти выражения подставить в уравнения (11), учесть, что 


$\frac{\partial a_{\alpha \beta }}{\partial y^\gamma }=\Gamma_{\alpha


\gamma \beta }+\Gamma_{\beta \gamma \alpha }$, то после


несложных, но достаточно громоздких подсчетов, получим


\begin{equation}


\lambda_{i,jk}=-\lambda_sR_{kij}^s+\psi_jg_{ik}+\psi_kg_{ij}+(\mu \Omega


_{ij})_{,k}+(\mu \Omega_{ik})_{,j}-(\mu \Omega_{jk})_{,i}. 


\end{equation}





\begin{thm}


Для того чтобы риманово пространство $V_n\subset V_{n+1}$


допускало бесконечно малые деформации, необходимо и достаточно,


чтобы в нем существовали вектор $\lambda_i$ и инвариант $\mu$,


удовлетворяющие уравнениям $(13)$.


\end{thm}





Система (13) есть координатная запись (11) в базисе 


$\{y_{,i}^\alpha ,\eta^\alpha \}$, а потому она дает необходимые


и достаточные условия геодезической деформации гиперповерхности


$V_n$. Будем называть систему (13) {\it основной} \cite{ga89}.





Симметрированием (13) по индексам $i$ и $j$ получим выражение


\begin{equation}


(\lambda_{i,j}+\lambda_{j,i}-2\mu \Omega_{ij})_{,k}=2\psi_kg_{ij}+\psi


_ig_{jk}+\psi_jg_{ik}.


\end{equation}


Аналогичными подсчетами как для уравнения (10), можно доказать, что


системы (13) и (14) эквивалентны. А (14) показывает, что для


тензора


$a_{ij}=\lambda_{i,j}+\lambda_{j,i}-2\mu \Omega_{ij}-2\psi g_{ij}$


она превращается в (9), что еще раз подтверждает теорему 2.





Для тензора


$h_{ij}=\lambda_{i,j}+\lambda_{j,i}-2\mu \Omega_{ij}$


уравнения (14) запишутся в форме (8).


В заключение этого пункта сделаем замечание.





Если $V_n$ совпадает с $V_m$, то (3) определяют бесконечно


малые преобразования $V_n$, изучавшиеся ранее \cite{ei}, \cite{kaj}.


В этом случае $y^\alpha =x^\alpha $, в (12) \ $\mu =0$, 


$\xi^\alpha =\lambda^\alpha $, а из (13) получается хорошо


известная система для бесконечно малых преобразований, которые


сохраняют геодезические


$\xi_{i,jk}=-\xi_sR_{kij}^s+\psi_ig_{jk}+\psi_jg_{ik}$. 


\medskip





{\bf 6. Система уравнений типа Коши для геодезических деформаций


гиперповерхности.} 


Вопрос о бесконечно малых геодезических деформациях


гиперповерхности риманова пространства сводится к решению


линейной системы (13), правая часть которых содержит, кроме


$\lambda_i$, производные от неизвестных $\mu$ и $\psi_i$. Эта


система допускает сведение к смешанной тензорной системе типа


Коши от $n^2+3n+2$ неизвестных функций.





Получению этой системы и посвящена эта часть.


В \cite{ga89} такая система получена для случая


невырожденного второго основного тензора гиперповерхности


$\Omega_{ij}$. Здесь покажем, что такая система существует при


более слабых условиях. В дальнейшем будем предполагать, что имеет место


$$


\rank\|\Omega_{ij}\|>2.


$$





Предварительно введем 


\begin{equation}


\lambda_{ij}=\lambda_{i,j},\quad


\mu_i=\mu_{,i}.


\end{equation}





В силу (15) формулы (13) перепишутся в следующем виде:


\begin{equation}


\lambda_{ij,k}=-\lambda_sR_{kij}^s+g_{ij}\psi_k+g_{ik}\psi_j+


\mu (\Omega_{ij,k}+\Omega_{ik,j}-\Omega_{jk,i})+


\mu_k\Omega_{ij}+\mu_j\Omega_{ik}-\mu_i\Omega_{jk}.


\end{equation}


Условия интегрируемости уравнений (16) имеют вид


\begin{gather}


\lambda_{sj}R^s_{ikl}+\lambda_{si}R^s_{jkl}=


\lambda{}_{s[k}R{}^s_{l]ij}+


\lambda_s(R^s_{lij,k}-R^s_{kij,l})+\notag \\


+g_{ik}\psi_{j,l}-g_{il}\psi_{j,k}+


\mu\Omega_{s(i} R^s_{j)kl}+


(\mu\Omega_{ik})_{,jl}+(\mu\Omega_{jl})_{,ik}-


(\mu\Omega_{il})_{,jk}-(\mu\Omega_{jk})_{,il},


\end{gather}


где круглыми и квадратными скобками обозначаем симметрирование и


альтернирование соответственно.





Если (17) просимметрируем по индексам $i$ и $j$, а затем свернем


с $g^{il}$ ($\|g^{ij}\|=\|g_{ij}\|^{-1}$), то получим


\begin{equation}


n\psi_{j,k}=h_{js}R^s_k-h_{ps}R^s{}_{jk}{}^p+\omega g_{jk},


\end{equation}


где $R^s_k=R^s_{ijk}g^{ij}$, $\omega$ --- некоторый инвариант,


индексы подняты тензором $g^{ij}$, 


$h_{ij}=\lambda_{(ij)}-2\mu\,\Omega_{ij}$. 





Найдем условия на неизвестный инвариант $\omega$.


Условия интегрируемости (18) приводятся к виду


$$


(n+3)\psi_sR_{jkl}^s=


\psi_sR{}_{[k}^sg{}_{l]j}+h_{js}R_{[k,l]}^s-h_{sp}R{}^s


_{\cdot j[k}{}^p_{\cdot,l]}+g_{j[k}\omega_{,l]}. 


$$


Сворачивая последнее выражение с $g^{jk}$, получим


\begin{equation}


(n-1)\omega_{,l}=2(n+1)\psi_sR_l^s+h_{js}R{}^s_{\cdot l,}{}^j_{\cdot}


-h_{sp}R{}^{sk}_{\cdot \cdot}{}_l{}^p_{\cdot, k}. 


\end{equation}


Если при помощи (18) исключить $\psi_{j,k}$ из (17), то получим


\begin{equation}


\lambda_{sp}A_{ijkl}^{sp}+\lambda_sB_{ijkl}^s+


\mu_sC_{ijkl}^s+\omega D_{ijkl}+\mu E_{ijkl}= 


\mu_{j,l}\Omega_{ik}-\mu_{j,k}\Omega_{il}-


\mu_{i,l}\Omega_{jk}+\mu_{i,k}\Omega_{jl}, 


\end{equation}


где 


\begin{gather*}


A_{ijkl}^{sp}\overset{\mathrm{def}}{=}{}


\delta_j^pR_{ikl}^s+\delta_i^sR_{jkl}^p+\delta


_l^pR_{kij}^s+\delta_k^sR_{lji}^p-\frac {1}{n}g_{ik}


(\delta{}_j^{(s}R{}_l^{p)}-R{}_{\cdot j}^{(s}{}_{l \cdot}^{p)})


+\frac 1ng_{il}


(\delta{}_j^{(s}R{}_k^{p)}-R{}^{(s}_{\cdot j k}{}^{p)}_\cdot),\\


%


B_{ijkl}^s\overset{\mathrm{def}}{=}{} R_{kij,l}^s-R_{lij,k}^s,\quad


C_{ijkl}^s\overset{\mathrm{def}}{=}{}


\delta{}_{[i}^s\Omega{}_{j]\,[k,l]}+\delta{}_{[k}^s\Omega{}


_{l]\,[i,j]}, \quad


D_{ijkl}\overset{\mathrm{def}}{=}{}-\frac 1n(g_{ik}g_{jl}-g_{il}g_{jk}),\\


%


E_{ijkl}\overset{\mathrm{def}}{=}{}-\Omega{}_{s(i}R{}_{j)kl}^s-


\Omega{}_{k[i,j]l}+\Omega{}_{l[i,j]k}+


\frac 2n\Omega{}_{js}g{}_{i[k}R{}_{l]}^s+


\frac 2n\Omega{}_{sp}


R{}^s_{\cdot j]k}{}^p_{\cdot}g{}_{l]i}.


\end{gather*}


Существенно отметить, что пять последних тензоров в конечном


счете определяются либо только метрикой $V_n$, либо метрикой и 


$\Omega_{ij}$.








Ввиду того, что $\rank\|\Omega_{ij}\|>2$,


в $V_n$, очевидно, существуют хотя бы три некомпланарные


векторные поля $a^i$, $b^i$, $c^i$ такие, что


$$


\Omega_{ij}a^ia^j=e_a=\pm 1,\quad 


\Omega_{ij}b^ib^j=e_b=\pm 1,\quad


\Omega_{ij}c^ic^j=e_c=\pm 1,\quad


\Omega_{ij}a^ib^j=\Omega_{ij}a^ic^j=\Omega_{ij}b^ic^j=0.


$$


Заметим, что эти векторы определяются $\Omega_{ij}$. Несмотря на


то, что они определены неоднозначно, их вправе считать


объектами пространства $V_n$, т.\,е. определенными через метрику и


вторую форму $V_n$.





Поочередно умножим (20) на $a^ia^k$, $b^ib^k$, $c^ic^k$, а затем


просуммируем по $i$ и~$k$. Получим


\begin{align}


e_a\mu_{j,l}&=\up{a}{\mu}_j\up{a}{\Omega}_l+\up{a}{\mu}_l\up{a}{\Omega}_j-


\up{a}{\mu}\Omega_{jl}+ 


\up{1}{T}_{jl}(\lambda_{sp},\lambda_s,\mu_s,\mu ,\omega),\tag{21\text{a}}\\


%


e_b\mu_{j,l}&=\up{b}{\mu}_j\up{b}{\Omega}_l+\up{b}{\mu}_l\up{b}{\Omega}_j-


\up{b}{\mu}\Omega_{jl}+ 


\up{2}{T}_{jl}(\lambda_{sp},\lambda_s,\mu_s,\mu ,\omega),\tag{21\text{b}}\\


%


e_c\mu_{j,l}&=\up{c}{\mu}_j\up{c}{\Omega}_l+\up{c}{\mu}_l\up{c}{\Omega}_j-


\up{c}{\mu}\Omega_{jl}+ 


\up{3}{T}_{jl}(\lambda_{sp},\lambda_s,\mu_s,\mu ,\omega ),\tag{21\text{c}}


\end{align}


где 


$\up{a}{\mu}_j\overset{\mathrm{def}}{=}{}\mu_{j,k}a^k$;


$\up{b}{\mu}_j\overset{\mathrm{def}}{=}{}\mu_{j,k}b^k$; 


$\up{c}{\mu}_j\overset{\mathrm{def}}{=}{}\mu_{j,k}c^k$; 


$\up{a}{\mu}\overset{\mathrm{def}}{=}{}\mu_{j,k}a^ja^k$; 


$\up{b}{\mu}\overset{\mathrm{def}}{=}{}\mu_{j,k}b^jb^k$; 


$\up{c}{\mu}\overset{\mathrm{def}}{=}{}\mu_{j,k}c^jc^k$;


$\up{a}{\Omega}_l\overset{\mathrm{def}}{=}{}\Omega_{lk}a^k$;


$\up{b}{\Omega}_l\overset{\mathrm{def}}{=}{}\Omega_{lk}b^k$; 


$\up{c}{\Omega}_l\overset{\mathrm{def}}{=}{}\Omega_{lk}c^k$; 


$\up{\sigma}{T}$ ($\sigma =1,2,\dotsc $) --- тензоры,


линейно выражающиеся через те объекты, которые указаны в качестве


аргумента, с коэффициентами, определенными метрикой $g_{ij}$


пространства $V_n$ и его второй квадратичной формой 


$\Omega_{ij}$, а также векторами $a^i$, $b^i$, $c^i$, которые


определены по существу $\Omega_{ij}$.


\addtocounter{equation}{1}





При помощи (21b) исключим в (21a) тензор $\mu_{j,l}$


\begin{equation}


e_a(\up{a}{\mu}_j\up{a}{\Omega}_l+\up{a}{\mu}_l\up{a}{\Omega}_j)-


e_b(\up{b}{\mu}_j\up{b}{\Omega}_l+\up{b}{\mu}_l\up{b}{\Omega}_j) -


(e_a\up{a}{\mu}-e_b\up{b}{\mu}) \Omega_{jl}+


e_a\up{1}{T}_{jl}-e_b\up{2}{T}_{jl}=0.


\end{equation}


Свернем (22) с $c^jc^l$, получим


$(e_a\up{a}{\mu}-e_b\up{b}{\mu})=e_c 


(e_a\up{1}{T}_{jl}-e_b\up{2}{T}_{jl})c^jc^l$.


Тогда после свертывания (22) с $a^ja^l$ будем иметь


$\up{a}{\mu}=\up{4}{T}(\lambda_{sp},\lambda_s,\mu_s,\mu ,\omega )$,


а в силу предыдущего и


$\up{b}{\mu}=\up{5}{T}(\lambda_{sp},\lambda_s,\mu_s,\mu ,\omega )$.





В итоге свертыванием (22) с $a^l$ убедимся, что 


\begin{equation}


\up{a}{\mu}_j=\alpha \up{b}{\Omega}_j+\up{6}{T}_j(\lambda_{sp},\lambda_s,\mu


_s,\mu ,\omega ),\tag{23\text{a}}


\end{equation}


где $\alpha$ --- некоторый инвариант.


Аналогичным путем анализом (21a) и (21c) получим 


\begin{equation}


\up{a}{\mu}_j=\beta \up{c}{\Omega}_j+\up{7}{T}_j(\lambda_{sp},\lambda_s,\mu


_s,\mu ,\omega ).\tag{23\text{b}}


\end{equation}


\addtocounter{equation}{1}








\noindent{}Вычитая (23b) из (23a), получим


$\alpha \up{b}{\Omega}_j-\beta \up{c}{\Omega}_j+\up{6}{T}_j-\up{7}{T}_j=0$.


Свертывая последнее с $b^j$, убедимся, что


$\alpha =\up{8}{T}(\lambda_{sp},\lambda_s,\mu_s,\mu ,\omega )$.


В итоге легко убедиться, что формула (21a) принимает вид:


\begin{equation}


\mu_{j,l}=\up{9}{T}_{jl}(\lambda_{sp},\lambda_s,\mu


_s,\mu ,\omega ).


\end{equation}





Совокупность уравнений (15), (16), (18), (19) и (24), которую 


кратко обозначим через $(A)$, дает линейное выражение первых


ковариантных производных от 


$\lambda_{ij}$, $\lambda_i$, $\psi_i$, $\mu_i$, $\mu$, $\omega$ 


через эти же тензоры с коэффициентами, однозначно определенными


$V_n$ и $V_{n+1}$, т.\,е. мы имеем систему типа Коши. Заметим, что


система $(A)$ носит инвариантный характер относительно выбора


координат в $V_n$. Итак, получена 





\begin{thm}


Для того чтобы риманово пространство $V_n\subset V_{n+1}$,


$\rank\|\Omega_{ij}\|>2$, 


допускало нетривиальные бесконечно малые геодезические


деформации, необходимо и достаточно, чтобы система уравнений $(A)$


имела нетривиальное решение.


\end{thm}





Условия интегрируемости всех уравнений системы $(A)$ суть линейные


однородные алгебраические уравнения относительно


$\lambda_{ij}$, $\lambda_i$, $\psi_i$, $\mu_i$, $\mu $, $\omega $. 


Обозначим их через $(B)$.


Такими же будут их дифференциальные продолжения любого порядка,


т.\,к. $(A)$ линейна. Продолжения обозначим


$(B_1)$, $(B_2)$ и т.\,д.








По аналогии с соответствующим результатом из теории геодезических


отображений \cite{si} получается 





\begin{thm}


Для того чтобы риманово пространство $V_n\subset V_{n+1}$ 


$(\rank\|\Omega_{ij}\|>2)$


допускало нетривиальные бесконечно малые геодезические деформации,


необходимо и достаточно, чтобы система линейных однородных


алгебраических уравнений


$(B),(B_1),(B_2),\dotsc,(B_s)$ $(s<n^2+3n+2)$


имела решение относительно


$\lambda_{ij}$, $\lambda_i$, $\psi_i$ $(\not\equiv0)$, $\mu_i$, $\mu$,


$\omega$.


\end{thm}





Отметим, что число $n^2+3n+2$ существенно понижается в случае,


когда $V_n$ не имеет постоянную кривизну. Об этом свидетельствуют


исследования \cite{si}, \cite{mige}, \cite{ya}, \cite{mikio}.








\begin{thebibliography}{99}





\bibitem{le}


Levi Civita T.


{\it Sulle transformationi delle equazioni dinamiche} //


Ann. Math. Milano. -- 1986. -- Ser.\,2. -- V.\,24. -- P.\,255--300.





\bibitem{ei}


Эйзенхарт Л.П. {\it Риманова геометрия}. -- М.: Ин. лит., 1948.





\bibitem{shi}


Широков П.А. {\it Избранные работы по геометрии}. -- Казань: Изд-во


Казанск. ун-та, 1966. -- 432~с.





\bibitem{no}


Норден А.П. {\it Пространства аффинной связности}. -- М.: Наука,


1976. -- 432~с.





\bibitem{pe}


Петров А.З. {\it Новые методы в общей теории относительности}. 


-- М.: Наука, 1966. -- 495~с. 





\bibitem{si}


Синюков Н.С. {\it Геодезические отображения римановых


пространств}. -- М.: Наука, 1979. -- 256~с.





\bibitem{mige}


Mike\v s J. {\it Geodesic mappings of affine-connected and


Riemannian spaces} // 


J. Math. Sci. New York. -- 1996. -- V.\,78. -- \No\,2. -- P.\,311--333.





\bibitem{ramiga}


Радулович Ж., Микеш Й., Гаврильченко М.Л. 


{\it Геодезические отображения и


деформации римановых пространств}. -- 


Izd. CID, Podgorica, Одесса: Изд-во Одесск. ун-та, 1997. -- 127~с.





\bibitem{siga}


Синюков Н.С., Гаврильченко М.Л.


{\it Бесконечно малые геодезические деформации поверхностей} //


Третья респ. конференция математиков Белоруссии, Минск, 1971.





\bibitem{ga72}


Гаврильченко М.Л.


{\it О геодезических деформациях гиперповерхностей} //


Тез. докл. V Всесоюзной конференции по современным проблемам


геометрии, Самарканд, 1972.





\bibitem{ga89}


Gavrilchenko M.L.


{\it Geodesic deformations of Riemannian space}  //


Diff. Geom. and its Appl. World Sci., Singapore, 1989.


-- P.\,47--53. 





\bibitem{ve}


Векуа И.Н. {\it Обобщенные аналитические функции}. --


М.: Физматгиз, 1959. -- 628 с.





\bibitem{kaj}


Кайгородов В.Р.


{\it О римановых пространствах $\up{\ast}{K}^s_n$} //


Тр. геометрич. семин. ВИНИТИ. -- 1974. -- T.\,5. -- C.\,359--373.





\bibitem{ya}


Yano K. {\it The theory of Lie derivatives and its applications}. -- 


Amst. North Holland publ. Groningen, Noordhoff, 1957. -- 299 p.





\bibitem{mikio}


Микеш Й., Киосак В.А. 


{\it О степени подвижности римановых пространств


относительно геодезических отображений} //


Геометрия погружен. многообразий. -- М.: МГПИ, 1986. -- C.\,35-39.





\bibitem{kid}


Kiosak V.A. 


{\it Geodetick\'a zobrazen\'\i\ Riemannov\'ych prostoru}:


Дис. Olomouc, 2002.





\end{thebibliography}





\vskip 2cm





\post{\it Одесский государственный}{\it Поступила}


\post{\it университет $($Украина, г.\,Одесса$)$}{25.03.2004}


\smallskip


\post{\it Одесская государственная академия}{}


\post{\it строительства и архитектуры}{}


\post{\it $($Украина, г.\,Одесса$)$}{}


\smallskip


\post{\it Оломоуцкий университет}{}


\post{\it $($Чешская республика$)$}{}





\label{end}


\end{document}





\thanks{ {\it Authors' address:\/}


1. Кафедра геометрии и топологии Одесского государственного


университета им. И.И. Мечникова,


 Дворянская, 2, Одесса, Украина;


2. Кафедра математики Одесской государственная академия


строительства и архитектуры,


 Ушинского, 2, Одесса, Украина;


3. Dept. of Algebra and Geometry, Fac. Sci., Palacky Univ., Tomkova 40,


 779 00 Olomouc, Czech Rep., e-mail: Mikes@inf.upol.cz.








